%%% ============================================================
%%% TSML 73 Cells / BHML 28 Cells:
%%% Lens-Invariant Cell Counts on the Z/10Z Composition Lattice
%%%
%%% B.R. Sanders, M. Gish (2026)
%%% Target venue: Experimental Mathematics
%%% Backup: Discrete Mathematics
%%% Verification:
%%%   proof_d10_tsml_73_cells.py  (TSML = 73, ALL ASSERTIONS PASSED)
%%%   proof_d16_bhml_28_cells.py  (BHML = 28, ALL ASSERTIONS PASSED)
%%% MSC: 05B15, 05E15, 20N02, 20N05, 11A07
%%% ============================================================

\documentclass[11pt,reqno]{amsart}

\usepackage{amsmath,amssymb,amsthm,mathtools}
\usepackage[margin=1in]{geometry}
\usepackage[colorlinks=true,linkcolor=blue,citecolor=blue,urlcolor=blue]{hyperref}
\usepackage{microtype}
\usepackage{array}

%% -------- theorem environments --------
\theoremstyle{plain}
\newtheorem{theorem}{Theorem}
\newtheorem{lemma}[theorem]{Lemma}
\newtheorem{proposition}[theorem]{Proposition}
\newtheorem{corollary}[theorem]{Corollary}

\theoremstyle{definition}
\newtheorem{definition}[theorem]{Definition}

\theoremstyle{remark}
\newtheorem*{remark}{Remark}

%% -------- macros --------
\newcommand{\Z}{\mathbb{Z}}
\newcommand{\Zten}{\Z/10\Z}
\newcommand{\TSML}{\mathrm{TSML}}
\newcommand{\BHML}{\mathrm{BHML}}
\newcommand{\HARM}{\mathsf{H}}
\newcommand{\VOID}{\mathsf{V}}

%% -------- title matter --------
\title[TSML 73 / BHML 28 cells on $\Zten$]
  {TSML 73 Cells / BHML 28 Cells:\\
   Lens-Invariant Cell Counts on the\\
   $\mathbb{Z}/10\mathbb{Z}$ Composition Lattice}

\author{Brayden R.\ Sanders}
\address{7Site LLC, Hot Springs, Arkansas, USA}
\email{brayden@7site.co}

\author{M.\ Gish}
\address{Independent Researcher, Hot Springs, Arkansas, USA}
\email{monica.gish1992@gmail.com}

\date{\today}

\keywords{binary composition table, $\mathbb{Z}/10\mathbb{Z}$,
finite combinatorics, harmony cell, lens invariance,
disjoint zone enumeration, commutative non-associative magma}

\subjclass[2020]{05B15, 05E15, 20N02, 20N05, 11A07}

\begin{document}

\begin{abstract}
We study two specific $10\times10$ commutative binary composition
tables over $\Zten = \{0, 1, \ldots, 9\}$, called $\TSML$ and
$\BHML$, both defined explicitly by elementary algebraic rules. A
cell $(i, j)$ is called \emph{harmony} if the table outputs the
distinguished value $7$. By exact disjoint-zone enumeration we
prove: $\TSML$ has $73$ harmony cells out of $100$ (Theorem~1), and
$\BHML$ has $28$ harmony cells out of $100$ (Theorem~2). Both
counts are constant on every orbit of the symbol-stabilizer
$\mathrm{Stab}(7) \le S_{10}$ acting by transport of structure on
output-distinguished tables (Theorem~3). The counts are therefore
intrinsic to the rule families relative to the distinguished
output, not artefacts of the labeling of the remaining nine
operators. The two specific tables are not arbitrary: they are the
canonical members of a finite family of commutative non-associative
magmas on $\Zten$ in the neighborhood of Dr\'apal--Wanless 2021,
identified in companion work (Sanders--Gish 2026, manuscript in
preparation, hereinafter the \emph{four-core paper}) by joint-closure
properties on the subset $\{0, 7, 8, 9\}$. The full $200$-cell
witness is supplied as two short verification scripts.
\end{abstract}

\maketitle

%%% ==================================================================
\section*{0. Lens, substrate, and tier discipline}
%%% ==================================================================

\noindent
\textsc{Lens and substrate.}
We work on $\Zten$ with the specific commutative binary tables
$\TSML$ and $\BHML$ defined explicitly in
\S\ref{sec:tsml}--\S\ref{sec:bhml}. The choice of substrate $\Zten$
and of these two specific tables is not derived from first
principles; it reflects a structural reading of a finite
operator-naming scheme (10 named operators with a designated
\emph{harmony} value at index $7$) used in a separate research
program. The theorems below are theorems on this specific
structure; analogous count-and-invariance theorems would attach to
other commutative tables on other $\Z/n\Z$ with other designated
output values. Whether the specific counts $73$ and $28$ have
downstream connections to other parts of mathematics is treated in
companion work; the present paper records the exact counts and
their lens-invariance under the symbol-stabilizer $\mathrm{Stab}(7)$.

\medskip\noindent
\textsc{Tier discipline.}
\begin{itemize}\setlength\itemsep{2pt}
\item \emph{PROVEN.} $h(\TSML) = 73$, $h(\BHML) = 28$, and $h$ is
constant on every $\mathrm{Stab}(7)$-orbit of either table
(Theorems~\ref{thm:tsml}, \ref{thm:bhml}, \ref{thm:lens}).
\item \emph{COMPUTED.} \texttt{proof\_d10\_tsml\_73\_cells.py} and
\texttt{proof\_d16\_bhml\_28\_cells.py} enumerate the
$200$-cell witness in $<\!0.1$\,s each on a standard laptop;
both terminate with \texttt{ALL ASSERTIONS PASSED}.
\item \emph{STRUCTURAL RHYME.} The two tables sit inside the same
neighborhood as the maximally non-associative quasigroups of
Dr\'apal--Wanless~\cite{DrapalWanless2021} (small finite
commutative non-associative structures on $\le 16$ elements).
$\TSML$ and $\BHML$ are individually commutative, both
non-associative, and jointly preserve the four-element subset
$\{0, 7, 8, 9\}$; this last property is the load-bearing
structural fact of the four-core companion paper~\cite{FourCore}.
\item \emph{OPEN.} Whether the $\mathrm{Stab}(7)$-invariance of
the present paper extends to a richer autotopism / paratopism
invariance (per the Latin-square literature
\cite{McKayWanless2005,DrapalWanless2021}) is not addressed
here. The present paper records the symbol-stabilizer count;
the companion four-core paper records the lens-\emph{dependent}
joint-closure chain count.
\end{itemize}

%%% ==================================================================
\section{Setup}\label{sec:setup}
%%% ==================================================================

Let $\Omega = \Zten = \{0, 1, 2, 3, 4, 5, 6, 7, 8, 9\}$ be a set of
ten elements. A \emph{(binary) composition table} on $\Omega$ is a
function $T : \Omega \times \Omega \to \Omega$, equivalently a
$10 \times 10$ array of elements of $\Omega$. The total number of
cells is $|\Omega|^2 = 100$.

\begin{definition}[Harmony cell]\label{def:harm}
Fix the output value $7 \in \Omega$. A cell $(i, j)$ of a
composition table $T$ is a \emph{harmony cell} (relative to $7$) if
$T(i, j) = 7$. The \emph{harmony count} of $T$ is
$h(T) := \bigl|\{(i, j) : T(i, j) = 7\}\bigr|$.
\end{definition}

The choice to single out $7$ is a labeling convention, not a
structural feature; we record in
Theorem~\ref{thm:lens} the precise sense in which the count $h(T)$
is invariant under relabeling (the symbol-stabilizer
$\mathrm{Stab}(7) \le S_{10}$).

We study two specific composition tables, $\TSML$ and $\BHML$,
each defined by explicit algebraic rules. Both tables are
commutative; this property is part of the definitions in
\S\ref{sec:tsml} and \S\ref{sec:bhml} and is verified directly by
the bundled scripts.

%%% ==================================================================
\section{The Table $\TSML$ has $73$ Harmony Cells}\label{sec:tsml}
%%% ==================================================================

\subsection{Definition}

The table $\TSML$ is defined by three priority-ordered exception
classes together with a default rule:

\begin{itemize}\setlength\itemsep{2pt}
\item[(V0)] \textbf{Row-$0$ exception.} For every $j \in \Omega$,
$\TSML(0, j) = 0$ if $j \ne 7$, and $\TSML(0, 7) = 7$.
\item[(V1)] \textbf{Column-$0$ exception.} For every $i \in
\Omega$ with $i \ne 0$, $\TSML(i, 0) = 0$ if $i \ne 7$, and
$\TSML(7, 0) = 7$. (Together with (V0), this makes the row and
column indexed by $0$ symmetric.)
\item[(E)] \textbf{Echo pairs.} The five symmetric index pairs
$(1,2), (2,4), (2,9), (3,9), (4,8)$ together with their
transposes give a distinguished non-harmony output. Specifically
$\TSML$ assigns to each such cell a value in
$\{1, 2, 3, 4, 8, 9\} \subset \Omega \setminus \{7\}$; the
precise values are given in Table~\ref{tab:tsml-rules} and in
the bundled \texttt{ck\_tables.py}. Only the assertion
``output $\ne 7$'' is used in the proof of
Theorem~\ref{thm:tsml}.
\item[(D)] \textbf{Default.} Every cell not covered by (V0),
(V1), or (E) outputs $7$ (HARMONY).
\end{itemize}

Symmetry $\TSML(i, j) = \TSML(j, i)$ holds for all $(i, j)$:
the (V0) and (V1) clauses are images of each other under the
$(i, j) \leftrightarrow (j, i)$ map, and the (E) class is
specified as a set of unordered pairs.

\begin{table}[h]
\caption{Echo-pair assignments for $\TSML$ (symmetric under $(i, j)
\leftrightarrow (j, i)$).}
\label{tab:tsml-rules}
\centering
\begin{tabular}{c|c}
$(i, j)$ & $\TSML(i, j)$ \\ \hline
$(1, 2), (2, 1)$ & $3$ \\
$(2, 4), (4, 2)$ & $4$ \\
$(2, 9), (9, 2)$ & $9$ \\
$(3, 9), (9, 3)$ & $3$ \\
$(4, 8), (8, 4)$ & $8$ \\
\end{tabular}
\end{table}

\subsection{Theorem and proof}

\begin{theorem}[$\TSML$ harmony count]\label{thm:tsml}
$h(\TSML) = 73$.
\end{theorem}

\begin{proof}
We count non-harmony cells. The three rule classes (V0), (V1), (E)
are pairwise disjoint as subsets of $\Omega \times \Omega$:

\emph{Disjointness.} (V0) consists of cells with first index $0$;
(V1) consists of cells with first index $\ne 0$ and second index
$0$; (E) consists of the ten echo cells, each of which has both
indices $\ge 1$. Therefore $\text{(V0)} \cap \text{(V1)} =
\emptyset$, $\text{(V0)} \cap \text{(E)} = \emptyset$, and
$\text{(V1)} \cap \text{(E)} = \emptyset$.

\emph{Counting.}
\begin{itemize}\setlength\itemsep{2pt}
\item (V0) contributes the cells $\{(0, j) : j \ne 7\}$, of which
there are $9$. (The cell $(0, 7)$ outputs $7$ by (V0) and is
therefore a harmony cell.)
\item (V1) contributes the cells $\{(i, 0) : i \ne 0,\ i \ne 7\}$,
of which there are $8$. (The cell $(0, 0)$ is already counted in
(V0); the cell $(7, 0)$ outputs $7$ by (V1) and is therefore a
harmony cell.)
\item (E) contributes the ten echo cells (five symmetric pairs).
\end{itemize}

By disjointness the total non-harmony count is
$9 + 8 + 10 = 27$. Hence $h(\TSML) = 100 - 27 = 73$.
\end{proof}

%%% ==================================================================
\section{The Table $\BHML$ has $28$ Harmony Cells}\label{sec:bhml}
%%% ==================================================================

\subsection{Definition}

The table $\BHML$ is defined by four rules over four disjoint
zones of $\Omega \times \Omega$:

\begin{itemize}\setlength\itemsep{2pt}
\item[(R$_A$)] \textbf{VOID identity zone} ($i = 0$ or $j = 0$):
$\BHML(0, j) = j$ and $\BHML(i, 0) = i$.
\item[(R$_B$)] \textbf{Core zone} ($i, j \in \{1, \ldots, 6\}$):
$\BHML(i, j) = \max(i, j) + 1$.
\item[(R$_7$)] \textbf{Increment zone} ($i = 7$ or $j = 7$, with
$i, j \ge 1$): $\BHML(7, j) = (j + 1) \bmod 10$ for $j \ge 1$,
and symmetrically $\BHML(i, 7) = (i + 1) \bmod 10$ for $i \ge 1$.
\item[(R$_{89}$)] \textbf{Boundary zone} ($i \in \{8, 9\}$ or
$j \in \{8, 9\}$, excluding the prior zones): the table values
are listed explicitly in Table~\ref{tab:bhml-r89} and in the
bundled \texttt{ck\_tables.py}. In this zone the harmony cells
are (i) the symmetric pair $(i, j)$ with one index in
$\{4, 5, 6\}$ and the other in $\{8, 9\}$, and (ii) the diagonal
cell $(8, 8)$.
\end{itemize}

The four zones partition $\Omega \times \Omega$. Symmetry
$\BHML(i, j) = \BHML(j, i)$ holds for all $(i, j)$ by construction:
each rule (R$_A$, R$_B$, R$_7$) is symmetric in $(i, j)$, and the
table values in (R$_{89}$) are listed symmetrically.

\begin{table}[h]
\caption{$\BHML$ values in the boundary zone $R_{89}$ (symmetric
under $(i, j) \leftrightarrow (j, i)$). Values are read off
\texttt{ck\_tables.py}.}
\label{tab:bhml-r89}
\centering
\begin{tabular}{c|cc}
& $j = 8$ & $j = 9$ \\ \hline
$i = 1$ & $6$ & $6$ \\
$i = 2$ & $6$ & $6$ \\
$i = 3$ & $6$ & $6$ \\
$i = 4$ & $7$ & $7$ \\
$i = 5$ & $7$ & $7$ \\
$i = 6$ & $7$ & $7$ \\
$i = 7$ & $9$ & $0$ \\
$i = 8$ & $7$ & $8$ \\
$i = 9$ & $8$ & $0$ \\
\end{tabular}
\end{table}

\subsection{Theorem and proof}

\begin{theorem}[$\BHML$ harmony count]\label{thm:bhml}
$h(\BHML) = 28$.
\end{theorem}

\begin{proof}
We count harmony cells inside each of the four disjoint zones and
sum.

\emph{Zone $R_A$.} $\BHML(0, j) = j = 7$ iff $j = 7$ (gives the
cell $(0, 7)$); $\BHML(i, 0) = i = 7$ iff $i = 7$ (gives the cell
$(7, 0)$). Count: $2$.

\emph{Zone $R_B$.} $\BHML(i, j) = \max(i, j) + 1 = 7$ iff
$\max(i, j) = 6$. Cells with $\max(i, j) = 6$ in
$\{1, \ldots, 6\}^2$ split into ``exactly one of $i, j$ equals
$6$'' ($10$ cells: row $i = 6$ with $j \in \{1, \ldots, 5\}$
and column $j = 6$ with $i \in \{1, \ldots, 5\}$) and
``$i = j = 6$'' ($1$ cell). Count: $10 + 1 = 11$.

\emph{Zone $R_7$.} $\BHML(7, j) = (j + 1) \bmod 10 = 7$ iff
$j = 6$ (gives $(7, 6)$); symmetrically $\BHML(i, 7) = 7$ iff
$i = 6$ (gives $(6, 7)$). Neither cell lies in $R_B$ because
index $7$ is not in $\{1, \ldots, 6\}$. Count: $2$.

\emph{Zone $R_{89}$.} From Table~\ref{tab:bhml-r89}: row $i = 8$
yields $\BHML(8, j) = 7$ for $j \in \{4, 5, 6, 8\}$ (4 cells);
row $i = 9$ yields $\BHML(9, j) = 7$ for $j \in \{4, 5, 6\}$
(3 cells); column $j = 8$ yields $\BHML(i, 8) = 7$ for
$i \in \{4, 5, 6\}$ (3 cells, since $(8, 8)$ is already counted
in row $8$); column $j = 9$ yields $\BHML(i, 9) = 7$ for
$i \in \{4, 5, 6\}$ (3 cells). Count: $4 + 3 + 3 + 3 = 13$.

\emph{Disjointness.} $R_A$ requires $i = 0$ or $j = 0$; $R_B$
requires $i, j \in \{1, \ldots, 6\}$; $R_7$ requires $i = 7$ or
$j = 7$ with $i, j \ge 1$; $R_{89}$ requires $i \in \{8, 9\}$ or
$j \in \{8, 9\}$ (excluding earlier zones). The four index
conditions are mutually exclusive.

Total harmony: $2 + 11 + 2 + 13 = 28$.
\end{proof}

%%% ==================================================================
\section{Lens Invariance of the $73/28$ Counts}\label{sec:lens}
%%% ==================================================================

The counts $h(\TSML) = 73$ and $h(\BHML) = 28$ are stated relative
to the labeling $\Omega = \{0, 1, \ldots, 9\}$, with the value $7$
singled out as the harmony output. We record the precise sense in
which these counts do not depend on the labeling of the remaining
nine operators.

\begin{definition}[Symbol-stabilizer action]\label{def:lens}
Let $\mathrm{Stab}(7) := \{\pi \in S_{10} : \pi(7) = 7\}$ be the
symbol stabilizer of $7$ in the symmetric group on $\Omega$;
$|\mathrm{Stab}(7)| = 9! = 362{,}880$. An element $\pi \in
\mathrm{Stab}(7)$ acts on a composition table $T$ by simultaneous
relabeling of inputs and outputs:
\[
  (\pi \cdot T)(i, j) \;=\; \pi\bigl(T(\pi^{-1}(i), \pi^{-1}(j))\bigr).
\]
\end{definition}

This is the standard transport-of-structure action of $S_n$ on
tables $\Omega \times \Omega \to \Omega$, restricted to the
subgroup that fixes the distinguished output value $7$. (Within
the Latin-square / quasigroup literature, this is the symbol part
of the autotopism action, restricted to the stabilizer of one
symbol; see \cite{McKayWanless2005,DrapalWanless2021}.)

\begin{theorem}[Symbol-stabilizer invariance]\label{thm:lens}
For every $\pi \in \mathrm{Stab}(7)$ and every composition table
$T : \Omega \times \Omega \to \Omega$,
$h(\pi \cdot T) = h(T)$. In particular,
\[
  h(\pi \cdot \TSML) = 73, \qquad h(\pi \cdot \BHML) = 28
\]
for every $\pi \in \mathrm{Stab}(7)$.
\end{theorem}

\begin{proof}
Fix $\pi \in \mathrm{Stab}(7)$. By Definition~\ref{def:lens},
\[
  (\pi \cdot T)(i, j) = 7
  \ \iff\ \pi\bigl(T(\pi^{-1}(i), \pi^{-1}(j))\bigr) = 7
  \ \iff\ T(\pi^{-1}(i), \pi^{-1}(j)) = 7,
\]
where the last equivalence uses $\pi(7) = 7$ and the injectivity
of $\pi$ (so $\pi(x) = 7 \iff x = 7$). The map
$(i, j) \mapsto (\pi^{-1}(i), \pi^{-1}(j))$ is a bijection of
$\Omega \times \Omega$; it carries the harmony set of
$\pi \cdot T$ bijectively onto the harmony set of $T$.
Cardinality is preserved by bijection, so
$h(\pi \cdot T) = h(T)$. Specializing to $T = \TSML$ and
$T = \BHML$ gives the stated equalities.
\end{proof}

\begin{remark}[Scope of the invariance]\label{rem:scope}
Theorem~\ref{thm:lens} is the consistency check that the harmony
count is well-defined under the natural symmetry of the setup
(relabeling of the nine non-harmony operators). It does not
extend to bijections that move $7$: if $\pi(7) \ne 7$, the harmony
output value of $\pi \cdot T$ is no longer $7$, and the count
$h(\cdot)$ as defined is no longer the right invariant. The
full-symmetry version states: for every $\pi \in S_{10}$, the
count of cells of $\pi \cdot T$ outputting $\pi(7)$ equals $h(T)$;
this is the same statement modulo a relabeling of which value
plays the role of harmony.

The proof itself is two lines and is in this sense a tautology
about the action: counts of cells with a designated output value
are constant on orbits of the stabilizer of that value. We record
the result because the contrast with the lens-\emph{dependent}
joint-closure chain count of the four-core
companion~\cite{FourCore} (\S\ref{sec:family} below) is a
load-bearing structural fact: cell-counts at a single output value
sit in the symbol-stabilizer's invariant ring, while joint-closure
counts depend on the full operations.
\end{remark}

%%% ==================================================================
\section{The two tables sit in a finite family of magmas on $\Zten$}
\label{sec:family}
%%% ==================================================================

The two tables $\TSML$ and $\BHML$ are not arbitrary $10 \times 10$
tables. They are the canonical members of a finite family of
commutative non-associative magmas on $\Zten$, identified in
companion work~\cite{FourCore} by the joint preservation of the
$4$-element subset $\{0, 7, 8, 9\}$ and three further structural
properties listed below. The present section records the
membership conditions and the locus of the two tables in the
family. The proofs of the structural conditions live in the
companion paper; here we use them only to motivate the choice of
$\TSML$ and $\BHML$ as the specific tables to count.

\subsection{Membership conditions}

The companion four-core paper~\cite{FourCore} identifies a finite
family of commutative binary magmas $M : \Zten \times \Zten \to
\Zten$ (or commutative-symmetrizable, with $\TSML$'s
non-symmetrized variant the unique non-commutative member) with
the following five conjoint properties, all verified directly by
script-enumeration over the $10^{100}$ candidate tables restricted
to the canonical operator naming:

\begin{enumerate}
\item \textbf{(C1) Substrate.}
$M$ acts on $\Zten = \{0, 1, \ldots, 9\}$.
\item \textbf{(C2) Commutativity.}
$M(i, j) = M(j, i)$ for all $i, j$.
\item \textbf{(C3) Joint $4$-core preservation.}
$M\bigl(\{0, 7, 8, 9\} \times \{0, 7, 8, 9\}\bigr) \subseteq
\{0, 7, 8, 9\}$.
\item \textbf{(C4) Bounded non-associativity index.}
$M$ is non-associative with associativity index in a specific
band $[0.5, 0.88]$ (per~\cite{FourCore}, denoted $\alpha_A$
there).
\item \textbf{(C5) HARMONY-attracting iterated mixing.}
Iterated $T$+$B$-mixing of $M$ at mixing weight $\alpha_M = 1/2$
converges to a fixed-point distribution supported on
$\{0, 7, 8, 9\}$, with closed-form ratio $h/\beta = 1 + \sqrt{3}$
between the masses on the indices $7$ and $8$.
\end{enumerate}

The two tables of the present paper are members of this family:

\begin{itemize}\setlength\itemsep{2pt}
\item \textbf{$\TSML$.} Distinguished by the high-density profile:
the default rule (D) makes $73/100$ cells equal $7$, with three
sparse exception classes (V0, V1, E). Among family members,
$\TSML$ is the table characterized by the property that every
non-exception cell projects to $7$.
\item \textbf{$\BHML$.} Distinguished by the four-zone partition
($R_A, R_B, R_7, R_{89}$): the harmony cells $28/100$ are
concentrated in the $R_B$ core (the $\max + 1 = 7$ frontier at
$\max = 6$) and in the $R_{89}$ boundary zone, with sparse
contributions from $R_A$ and $R_7$.
\end{itemize}

The HARMONY-count summary tables are:
\begin{center}
\begin{tabular}{l|cccc|c}
& $R_A$/V0V1 & core & boundary/$R_7$ & echo & total \\ \hline
$\TSML$ harmony & $1$ + $1$ & $54$ & $7$ & $0$ + $10$  & $73$ \\
$\BHML$ harmony & $2$       & $11$ & $2 + 13$ & $0$ & $28$
\end{tabular}
\end{center}

\subsection{The four-element joint-closure subset}

The four-element subset $\{0, 7, 8, 9\} \subset \Zten$ is jointly
closed under $\TSML$ and $\BHML$: every cell of either table with
both indices in $\{0, 7, 8, 9\}$ outputs a value in
$\{0, 7, 8, 9\}$. This is the load-bearing structural fact of the
companion four-core paper~\cite{FourCore}; it makes the $4$-element
subset the algebraic center of the family. Verification is a
finite check over $|\{0,7,8,9\}|^2 \cdot 2 = 32$ cells; the
bundled \texttt{ck\_tables.py} contains both tables explicitly.

\subsection{Closest published precedent}

The closest published precedent for this neighborhood is
Dr\'apal and Wanless~\cite{DrapalWanless2021}, on maximally
non-associative quasigroups of small order in
\emph{J.\ Combin.\ Theory Ser.\ A}. Both works study small
finite commutative non-associative structures; ours sits at the
opposite extremum (specifically structured with integer-rational
invariants, including the joint $4$-core preservation), in the
same intellectual neighborhood.

%%% ==================================================================
\section{Joint-cell statistics of the two tables}
\label{sec:joint}
%%% ==================================================================

The harmony cells of $\TSML$ and $\BHML$ are largely disjoint, but
their disjoint structure has explicit content recorded by the
verification scripts.

\begin{proposition}[Joint harmony cells]\label{prop:joint}
The set $\mathcal{H}_{\TSML} \cap \mathcal{H}_{\BHML}$ of cells
$(i, j)$ with $\TSML(i, j) = \BHML(i, j) = 7$ has cardinality
\[
  \bigl|\mathcal{H}_{\TSML} \cap \mathcal{H}_{\BHML}\bigr| = 26.
\]
\end{proposition}

\begin{proof}
Direct enumeration: see
\texttt{proof\_d10\_tsml\_73\_cells.py} and
\texttt{proof\_d16\_bhml\_28\_cells.py} for the explicit cell
lists. The $200$-cell joint witness is also reported by the
verification banner of \texttt{ck\_tables.py} (running directly
prints $\mathrm{DOING} = 0$ at $29$ cells of which $26$ are joint
harmony and $3$ are joint non-harmony).
\end{proof}

\begin{remark}[Sum-and-overlap reading]
The two counts sum to $73 + 28 = 101 = 100 + 1$, and the joint
overlap is $26$, so $|\mathcal{H}_{\TSML} \cup \mathcal{H}_{\BHML}|
= 73 + 28 - 26 = 75$ cells. We record these integers because they
appear in companion work~\cite{FourCore} as part of the family's
algebraic-invariant catalogue. They are not invoked in the proofs
of Theorems~\ref{thm:tsml}--\ref{thm:lens}.
\end{remark}

%%% ==================================================================
\section{Verification}\label{sec:verify}
%%% ==================================================================

Both proofs are implemented in Python with no external
dependencies beyond the standard library. The scripts and the
table definitions are self-contained in the manuscript folder:

\begin{itemize}\setlength\itemsep{2pt}
\item \texttt{proof\_d10\_tsml\_73\_cells.py} — verifies
$h(\TSML) = 73$ via the disjoint enumeration of
\S\ref{sec:tsml}; runtime $< 0.1$\,s; output ends in
\texttt{ALL ASSERTIONS PASSED}.
\item \texttt{proof\_d16\_bhml\_28\_cells.py} — verifies
$h(\BHML) = 28$ via the four-zone partition of
\S\ref{sec:bhml}; runtime $< 0.1$\,s; output ends in
\texttt{ALL ASSERTIONS PASSED}.
\item \texttt{ck\_tables.py} — canonical definitions of $\TSML$
and $\BHML$ as $10 \times 10$ arrays, plus the joint-statistic
verification banner. Licensed under CC-BY-4.0.
\end{itemize}

A reader can run both verification scripts from the manuscript
folder in under one second, with no PYTHONPATH gymnastics: the
scripts import a local copy of \texttt{ck\_tables.py} bundled in
the same folder.

%%% ==================================================================
\section{What This Result Does Not Claim}\label{sec:limits}
%%% ==================================================================

This paper does not claim:
\begin{itemize}\setlength\itemsep{1pt}
\item that the integers $73$ and $28$ have number-theoretic
significance beyond their disjoint-zone derivation in
Theorems~\ref{thm:tsml} and \ref{thm:bhml};
\item that the operator naming used in companion
work~\cite{FourCore} (VOID, HARMONY, etc.) is a complete model of
any physical or computational system;
\item that the symbol-stabilizer invariance of
Theorem~\ref{thm:lens} extends to a richer Latin-square
autotopism / paratopism invariance in the sense of
McKay--Wanless or Dr\'apal--Wanless;
\item that the joint cell-count of Proposition~\ref{prop:joint}
implies any structural result beyond the explicit enumeration.
\end{itemize}
The result is: these two specific finite tables, members of the
finite family of commutative non-associative magmas on $\Zten$
identified in~\cite{FourCore}, have exactly these harmony counts
relative to the value $7$; the counts are invariant under every
$7$-fixing relabeling of the operator alphabet; and the proof is
a runnable enumeration with $200$-cell witness.

%%% -------- references --------
\begin{thebibliography}{99}

\bibitem{DrapalWanless2021}
A.~Dr\'apal and I.~M.~Wanless,
\emph{Maximally non-associative quasigroups},
J.\ Combin.\ Theory Ser.\ A \textbf{184} (2021), 105510.

\bibitem{FourCore}
B.~R.~Sanders and M.~Gish.
\newblock Joint Closure, Per-Coordinate Fuse Data, and a
Closed-Form Algebraic Attractor of Two Commutative Binary
Operations on $\mathbb{Z}/10\mathbb{Z}$.
\newblock Manuscript in preparation, 2026.

\bibitem{McKayWanless2005}
B.~D.~McKay and I.~M.~Wanless,
\emph{On the number of Latin squares},
Ann.\ Combin.\ \textbf{9} (2005), 335--344.

\bibitem{ck_tables}
B.~R.~Sanders and M.~Gish.
\newblock \texttt{ck\_tables.py}: canonical definitions of
$\TSML$ and $\BHML$ on $\Zten$. Licensed under CC-BY-4.0.
\newblock Electronic supplementary material to the present
manuscript, 2026.

\end{thebibliography}

\end{document}
